\documentclass[a4paper,12pt,twoside]{memoir}

% Castellano
\usepackage[spanish,es-tabla]{babel}
\selectlanguage{spanish}
\usepackage[utf8]{inputenc}
\usepackage[T1]{fontenc}
\usepackage{lmodern} % scalable font
\usepackage{microtype}
\usepackage{placeins}

\RequirePackage{booktabs}
\RequirePackage[table]{xcolor}
\RequirePackage{xtab}
\RequirePackage{multirow}

% Landscape
\usepackage{pdflscape}

% Links
\usepackage[colorlinks]{hyperref}
\hypersetup{
	allcolors = {red}
}

% Ecuaciones
\usepackage{amsmath}

% Rutas de fichero / paquete
\newcommand{\ruta}[1]{{\sffamily #1}}

% Párrafos
\nonzeroparskip


% Imagenes
\usepackage{graphicx}
\newcommand{\imagen}[2]{
	\begin{figure}[!h]
		\centering
		\includegraphics[width=0.9\textwidth]{#1}
		\caption{#2}\label{fig:#1}
	\end{figure}
	\FloatBarrier
}

\newcommand{\imagenflotante}[2]{
	\begin{figure}%[!h]
		\centering
		\includegraphics[width=0.9\textwidth]{#1}
		\caption{#2}\label{fig:#1}
	\end{figure}
}



% El comando \figura nos permite insertar figuras comodamente, y utilizando
% siempre el mismo formato. Los parametros son:
% 1 -> Porcentaje del ancho de página que ocupará la figura (de 0 a 1)
% 2 --> Fichero de la imagen
% 3 --> Texto a pie de imagen
% 4 --> Etiqueta (label) para referencias
% 5 --> Opciones que queramos pasarle al \includegraphics
% 6 --> Opciones de posicionamiento a pasarle a \begin{figure}
\newcommand{\figuraConPosicion}[6]{%
  \setlength{\anchoFloat}{#1\textwidth}%
  \addtolength{\anchoFloat}{-4\fboxsep}%
  \setlength{\anchoFigura}{\anchoFloat}%
  \begin{figure}[#6]
    \begin{center}%
      \Ovalbox{%
        \begin{minipage}{\anchoFloat}%
          \begin{center}%
            \includegraphics[width=\anchoFigura,#5]{#2}%
            \caption{#3}%
            \label{#4}%
          \end{center}%
        \end{minipage}
      }%
    \end{center}%
  \end{figure}%
}

%
% Comando para incluir imágenes en formato apaisado (sin marco).
\newcommand{\figuraApaisadaSinMarco}[5]{%
  \begin{figure}%
    \begin{center}%
    \includegraphics[angle=90,height=#1\textheight,#5]{#2}%
    \caption{#3}%
    \label{#4}%
    \end{center}%
  \end{figure}%
}
% Para las tablas
\newcommand{\otoprule}{\midrule [\heavyrulewidth]}
%
% Nuevo comando para tablas pequeñas (menos de una página).
\newcommand{\tablaSmall}[5]{%
 \begin{table}
  \begin{center}
   \rowcolors {2}{gray!35}{}
   \begin{tabular}{#2}
    \toprule
    #4
    \otoprule
    #5
    \bottomrule
   \end{tabular}
   \caption{#1}
   \label{tabla:#3}
  \end{center}
 \end{table}
}

%
%Para el float H de tablaSmallSinColores
\usepackage{float}

%
% Nuevo comando para tablas pequeñas (menos de una página).
\newcommand{\tablaSmallSinColores}[5]{%
 \begin{table}[H]
  \begin{center}
   \begin{tabular}{#2}
    \toprule
    #4
    \otoprule
    #5
    \bottomrule
   \end{tabular}
   \caption{#1}
   \label{tabla:#3}
  \end{center}
 \end{table}
}

\newcommand{\tablaApaisadaSmall}[5]{%
\begin{landscape}
  \begin{table}
   \begin{center}
    \rowcolors {2}{gray!35}{}
    \begin{tabular}{#2}
     \toprule
     #4
     \otoprule
     #5
     \bottomrule
    \end{tabular}
    \caption{#1}
    \label{tabla:#3}
   \end{center}
  \end{table}
\end{landscape}
}

%
% Nuevo comando para tablas grandes con cabecera y filas alternas coloreadas en gris.
\newcommand{\tabla}[6]{%
  \begin{center}
    \tablefirsthead{
      \toprule
      #5
      \otoprule
    }
    \tablehead{
      \multicolumn{#3}{l}{\small\sl continúa desde la página anterior}\\
      \toprule
      #5
      \otoprule
    }
    \tabletail{
      \hline
      \multicolumn{#3}{r}{\small\sl continúa en la página siguiente}\\
    }
    \tablelasttail{
      \hline
    }
    \bottomcaption{#1}
    \rowcolors {2}{gray!35}{}
    \begin{xtabular}{#2}
      #6
      \bottomrule
    \end{xtabular}
    \label{tabla:#4}
  \end{center}
}

%
% Nuevo comando para tablas grandes con cabecera.
\newcommand{\tablaSinColores}[6]{%
  \begin{center}
    \tablefirsthead{
      \toprule
      #5
      \otoprule
    }
    \tablehead{
      \multicolumn{#3}{l}{\small\sl continúa desde la página anterior}\\
      \toprule
      #5
      \otoprule
    }
    \tabletail{
      \hline
      \multicolumn{#3}{r}{\small\sl continúa en la página siguiente}\\
    }
    \tablelasttail{
      \hline
    }
    \bottomcaption{#1}
    \begin{xtabular}{#2}
      #6
      \bottomrule
    \end{xtabular}
    \label{tabla:#4}
  \end{center}
}

%
% Nuevo comando para tablas grandes sin cabecera.
\newcommand{\tablaSinCabecera}[5]{%
  \begin{center}
    \tablefirsthead{
      \toprule
    }
    \tablehead{
      \multicolumn{#3}{l}{\small\sl continúa desde la página anterior}\\
      \hline
    }
    \tabletail{
      \hline
      \multicolumn{#3}{r}{\small\sl continúa en la página siguiente}\\
    }
    \tablelasttail{
      \hline
    }
    \bottomcaption{#1}
  \begin{xtabular}{#2}
    #5
   \bottomrule
  \end{xtabular}
  \label{tabla:#4}
  \end{center}
}



\definecolor{cgoLight}{HTML}{EEEEEE}
\definecolor{cgoExtralight}{HTML}{FFFFFF}

%
% Nuevo comando para tablas grandes sin cabecera.
\newcommand{\tablaSinCabeceraConBandas}[5]{%
  \begin{center}
    \tablefirsthead{
      \toprule
    }
    \tablehead{
      \multicolumn{#3}{l}{\small\sl continúa desde la página anterior}\\
      \hline
    }
    \tabletail{
      \hline
      \multicolumn{#3}{r}{\small\sl continúa en la página siguiente}\\
    }
    \tablelasttail{
      \hline
    }
    \bottomcaption{#1}
    \rowcolors[]{1}{cgoExtralight}{cgoLight}

  \begin{xtabular}{#2}
    #5
   \bottomrule
  \end{xtabular}
  \label{tabla:#4}
  \end{center}
}




\graphicspath{ {./img/} }

% Capítulos
\chapterstyle{bianchi}
\newcommand{\capitulo}[2]{
	\setcounter{chapter}{#1}
	\setcounter{section}{0}
	\chapter*{#2}
	\addcontentsline{toc}{chapter}{#2}
	\markboth{#2}{#2}
}

% Apéndices
\renewcommand{\appendixname}{Apéndice}
\renewcommand*\cftappendixname{\appendixname}

\newcommand{\apendice}[1]{
	%\renewcommand{\thechapter}{A}
	\chapter{#1}
}

\renewcommand*\cftappendixname{\appendixname\ }

% Formato de portada
\makeatletter
\usepackage{xcolor}
\newcommand{\tutor}[1]{\def\@tutor{#1}}
\newcommand{\course}[1]{\def\@course{#1}}
\definecolor{cpardoBox}{HTML}{E6E6FF}
\def\maketitle{
  \null
  \thispagestyle{empty}
  % Cabecera ----------------
\noindent\includegraphics[width=\textwidth]{cabecera}\vspace{1cm}%
  \vfill
  % Título proyecto y escudo informática ----------------
  \colorbox{cpardoBox}{%
    \begin{minipage}{.8\textwidth}
      \vspace{.5cm}\Large
      \begin{center}
      \textbf{TFG del Grado en Ingeniería Informática}\vspace{.6cm}\\
      \textbf{\LARGE\@title{}}
      \end{center}
      \vspace{.2cm}
    \end{minipage}

  }%
  \hfill\begin{minipage}{.20\textwidth}
    \includegraphics[width=\textwidth]{escudoInfor}
  \end{minipage}
  \vfill
  % Datos de alumno, curso y tutores ------------------
  \begin{center}%
  {%
    \noindent\LARGE
    Presentado por \@author{}\\ 
    en Universidad de Burgos --- \@date{}\\
    Tutor: \@tutor{}\\
  }%
  \end{center}%
  \null
  \cleardoublepage
  }
\makeatother


% Datos de portada
\title{Sistema de Recomendación basado en Aprendizaje Profundo \\Documentación Técnica}
\author{Raúl Negro Carpintero}
\tutor{Bruno Baruque Zanón}
\date{\today}

\begin{document}

\maketitle



\cleardoublepage



%%%%%%%%%%%%%%%%%%%%%%%%%%%%%%%%%%%%%%%%%%%%%%%%%%%%%%%%%%%%%%%%%%%%%%%%%%%%%%%%%%%%%%%%



\frontmatter


\clearpage

% Indices
\tableofcontents

\clearpage

\listoffigures

\clearpage

\listoftables

\clearpage

\mainmatter

\appendix

\apendice{Plan de Proyecto Software}

\section{Introducción}
La planificación es un punto fundamental en cualquier proyecto software. En ella, se tiene que estimar qué cosas hay que hacer y cuánto tiempo y dinero va a llevar terminarlas. De esta manera, se podrá comprobar si se llega a tiempo a los distintos plazos y así evitar retrasos, lo que supondría pérdidas de dinero y tiempo.

La planificación se puede dividir en dos partes clave:
\begin{itemize}
\tightlist
\item \textbf{Planificación temporal:} consiste en dividir el proyecto en etapas (\textit{sprints}) con una duración de x días. En estas etapas hay que dejar claro qué cosas se quieren conseguir y cuánto tiempo o esfuerzo se va a dedicar a cada una de ellas. Los \textit{sprints} no tienen por qué durar siempre lo mismo, se pueden ajustar a las necesidades de cada momento.
\item \textbf{Estudio de viabilidad:} con este estudio se podrá ver si el proyecto puede seguir adelante (es viable) o si no, con lo que habría que tomar medidas para cambiar la situación. Se divide en:
\begin{itemize}
\tightlist
\item \textbf{Viabilidad económica:} según la cual se deberá calcular por cuánto tiene que venderse el servicio para compensar el dinero, el tiempo y el esfuerzo puestos en el proyecto. Se incluye salario del personal y la compra de software de terceros necesario para la realización del proyecto.
\item \textbf{Viabilidad legal:} se estudia el licenciamiento del software que se va a utilizar para proceder dentro de la legalidad.
\end{itemize}
\end{itemize}

\section{Planificación temporal}
Para el desarrollo del proyecto se decidió seguir la metodología ágil \textit{Scrum}. Al ser un proyecto académico con una sola persona trabajando en él, no se siguen a raja tabla todos los rasgos característicos de \textit{Scrum}. Las características que se han seguido son:
\begin{itemize}
\tightlist
\item Desarrollo incremental en cada \textit{sprint}.
\item Cada \textit{sprint} tiene una duración aproximada de dos semanas.
\item Se realizan reuniones al final de cada \textit{sprint} con el objetivo de repasar lo hecho y pensar en los objetivos del siguiente \textit{sprint}.
\end{itemize}

A continuación se desarrolla el contenido de cada \textit{sprint}.

\subsection{Sprint 1 (- 13/12/2018)}\label{sprint-1}
Este \textit{sprint} se corresponde con el milestone \href{https://github.com/rnc0011/SistemaRecomendacionTFG/milestone/1}{Búsqueda de información inicial}. En él se estudia el capítulo dedicado a los sistemas de recomendación del libro \textit{Mining of Massive Datasets} \cite{miningDatasets}.

También se realiza el curso de \textit{fast.ai} \cite{fastai}. Además, se escoge la librería de \textit{LightFM} y se obtiene un modelo inicial.

\subsection{Sprint 2 (13/12/2018 - 18/01/2019)}\label{sprint-2}
Este \textit{sprint} se corresponde con el milestone \href{https://github.com/rnc0011/SistemaRecomendacionTFG/milestone/3}{Explorar modelo con LightFM}. En él se intenta comprender el funcionamiento de \textit{LightFM} y se utiliza el conjunto de datos de \textit{Movielens} con el modelo.

\subsection{Sprint 3 (18/01/2019 - 16/02/2019)}\label{sprint-3}
Este \textit{sprint} se corresponde con el milestone \href{https://github.com/rnc0011/SistemaRecomendacionTFG/milestone/4}{Recomendación híbrida y por contenido con LightFM}. En él se obtienen las primeras versiones del modelo híbrido y por contenido de \textit{LightFM}. 

También se obtiene un modelo inicial con \textit{PyTorch} y se empieza a pensar en las métricas que se usarán para evaluar los modelos.

\subsection{Sprint 4 (16/02/2019 - 03/03/2019)}\label{sprint-4}
Este \textit{sprint} se corresponde con el milestone \href{https://github.com/rnc0011/SistemaRecomendacionTFG/milestone/5}{Trabajar en la documentación}. En él se empieza a ampliar los apartados de la documentación relativos a los conjuntos de datos utilizados.

\subsection{Sprint 5 (03/03/2019 - 13/03/2019)}\label{sprint-5}
Este \textit{sprint} se corresponde con el milestone \href{https://github.com/rnc0011/SistemaRecomendacionTFG/milestone/6}{División de los datos en LightFM}. En él se dividen los conjuntos de datos en \textit{train} y \textit{test} y se obtienen de nuevo los modelos.

\subsection{Sprint 6 (13/03/2019 - 28/03/2019)}\label{sprint-6}
Este \textit{sprint} se corresponde con el milestone \href{https://github.com/rnc0011/SistemaRecomendacionTFG/milestone/7}{Medidas de calidad LightFM}. En él se obtienen las métricas finales para los modelos de \textit{LightFM}.

Además, se guardan las matrices y los modelos en archivos \textit{pickle} como medida de persistencia.

\subsection{Sprint 7 (28/03/2019 - 05/04/2019)}\label{sprint-7}
Este \textit{sprint} se corresponde con el milestone \href{https://github.com/rnc0011/SistemaRecomendacionTFG/milestone/8}{Últimos pasos con LightFM}. En él se intenta acabar todo lo relacionado con los modelos clásicos (tanto código como documentación).

\subsection{Sprint 8 (05/04/2019 - 24/04/2019)}\label{sprint-8}
Este \textit{sprint} se corresponde con el milestone \href{https://github.com/rnc0011/SistemaRecomendacionTFG/milestone/9}{Ampliar documentación}. En él se amplía la documentación relativa a \textit{LightFM} y se empiezan a obtener los primeros diagramas UML con la estructura del proyecto hasta el momento.

\subsection{Sprint 9 (24/04/2019 - 08/05/2019)}\label{sprint-9}
Este \textit{sprint} se corresponde con el milestone \href{https://github.com/rnc0011/SistemaRecomendacionTFG/milestone/10}{Cambios en el diseño, Flask y DL}. En él se empieza a obtener la interfaz web con \textit{Flask} y se escoge \textit{Spotlight} como la librería de aprendizaje profundo.

También se modifica la documentación para ir acorde a los cambios del código.

\subsection{Sprint 10 (08/05/2019 - 22/05/2019)}\label{sprint-10}
Este \textit{sprint} se corresponde con el milestone \href{https://github.com/rnc0011/SistemaRecomendacionTFG/milestone/11}{Primeros pasos con Spotlight}. En él se obtiene un primer modelo de \textit{Spotlight} y se unifican los archivos de Entrada y Persistencia para que solo haya uno de cada.

Se continúa trabajando en la documentación.

\subsection{Sprint 11 (22/05/2019 - 05/06/2019)}\label{sprint-11}
Este \textit{sprint} se corresponde con el milestone \href{https://github.com/rnc0011/SistemaRecomendacionTFG/milestone/12}{Completar modelos de Spotlight}. En él se obtienen el resto de modelos de \textit{Spotlight} y se generalizan los métodos de lectura de datos y obtención de matrices.

Se continúa trabajando en la documentación.

\subsection{Sprint 12 (05/06/2019 - 12/06/2019)}\label{sprint-12}
Este \textit{sprint} se corresponde con el milestone \href{https://github.com/rnc0011/SistemaRecomendacionTFG/milestone/13}{Diseño GUI}. En él se piensa qué forma va a tener la interfaz web obtenida a través de \textit{Flask}.

\subsection{Sprint 13 (12/06/2019 - 20/06/2019)}\label{sprint-13}
Este \textit{sprint} se corresponde con el milestone \href{https://github.com/rnc0011/SistemaRecomendacionTFG/milestone/14}{Obtención GUI}. En él se termina de crear la interfaz web y se empiezan a emplear los métodos obtenidos anteriormente para que la aplicación sea funcional.

\subsection{Sprint 14 (20/06/2019 - 03/07/2019)}\label{sprint-14}
Este \textit{sprint} se corresponde con el milestone \href{https://github.com/rnc0011/SistemaRecomendacionTFG/milestone/15}{Último milestone}. En él se termina todo lo que queda por hacer en el proyecto. La interfaz es totalmente funcional y la documentación está terminada.

\section{Estudio de viabilidad}
\subsection{Viabilidad económica}
En este apartado se emularán los costes que hubiera tenido una empresa si hubiera dedicado tiempo y recursos a desarrollar este proyecto, así como el beneficio que podría sacar gracias a su venta. 

Por lo tanto, los dos elementos fundamentales son: los costes y los beneficios.

\subsection{Costes}\label{costes}
Los costes se pueden dividir de la siguiente manera:

\textbf{Costes humanos:}

En este conjunto se incluye el sueldo de los trabajadores. Como el proyecto se ha podido llevar a cabo por una solo persona, se van a calcular los costes como si el equipo de trabajo estuviera compuesto por una persona. Además, la situación de esta persona sería la de recién graduado y sin experiencia laboral en este campo. La duración sería de unos 5 meses. 

Según un informe de \textit{Adecco} en 2015 \cite{sueldo-junior}, los informáticos junior cobraban alrededor de 17.248€ brutos al año. Sabiendo esto, el desglose quedaría:

\tablaSinColores{Coste de personal}{l l}{2}{coste-personal}
{\textbf{Concepto} & \textbf{Coste} \\}{
	Salario mensual neto & 662,61€ \\
	Seguridad Social (29,9\%) & 429,76€ \\ 
	Retención IRPF (24\%) & 344,96€ \\
	Salario mensual bruto & 1437.33€ \\
	\midrule
	Total 5 meses & 7.186,65€ \\
}

\subsection{Viabilidad legal}



\apendice{Especificación de Requisitos}

\section{Introducción}

\section{Objetivos generales}

\section{Catalogo de requisitos}

\section{Especificación de requisitos}



\apendice{Especificación de diseño}

\section{Introducción}
En este apéndice se explica cómo están conformados los datos que utilizan las librerías usadas en el proyecto, así como la forma en la que está estructurado el mismo.

\section{Diseño de datos}
Todos los datos que he utilizado a lo largo del proyecto están en formato .csv. Lo más normal es que para cada conjunto de datos tenga los siguientes archivos:
\begin{itemize}
\tightlist
\item \textit{ratings.csv} 
\item \textit{users.csv}
\item \textit{items.csv}
\end{itemize}
La estructura de estos archivos suele ser: \textit{idUser, idItem, rating, timestamp} para \textit{ratings.csv}, \textit{idUser, name, feature1, ..., featureN} para \textit{users.csv} y i\textit{dItem, name, feature1, ..., featureN} para \textit{items.csv}.

Para poder trabajar con los datos primero los paso a \textit{DataFrames} de \textit{pandas} \cite{dataframes}.

\subsection{Datos con LightFM}\label{datos-lightfm}
Una vez obtenidos los \textit{DataFrames} para cada \textit{.csv}, necesito convertirlos a \textit{Dataset} de \textit{LightFM} \cite{dataset-lightfm} para poder trabajar con ellos.

Esta clase se encarga de convertir los datos almacenados en los \textit{DataFrames} en \textit{matrices COO} y \textit{matrices CSR}.

\section{Diseño procedimental}

\section{Diseño arquitectónico}
Para la realización de este proyecto se ha seguido el patrón arquitectónico MVC (\textit{Modelo Vista Controlador}). El objetivo de este patrón es dividir el código en función de su propósito. Sus partes son:
\begin{itemize}
\tightlist
\item \textit{Modelo}: el acceso a los datos. Se corresponde con las clases de Entrada y Salida, que leen los datos para dárselo al sistema de recomendación y guardan los resultados.
\item \textit{Vista}: la visualización de los datos. Se corresponde con las clases de Interfaz, que muestran la información solicitada.
\item \textit{Controlador}: la manipulación de los datos. Se corresponde con los clases de Sistema, que crean los sistemas de recomendación gracias a los datos proporcionados por las clases de Entrada.
\end{itemize}

\imagen{patronMVC}{Esquema del patrón MVC}

La estructura del proyecto siguiendo este patrón quedaría de la siguiente forma:
\imagen{diagrama_src}{Diagrama UML del proyecto}
Por separado, los paquetes contienen:
\imagen{diagrama_vista}{Diagrama UML del paquete \textit{vista}}
\imagen{diagrama_controlador}{Diagrama UML del paquete \textit{controlador}}
\imagen{diagrama_modelo}{Diagrama UML del paquete \textit{modelo}}

\subsection{Diseño con LightFM}\label{diseño-lightfm}
Para la parte del modelo clásico, se tienen los siguientes archivos:
En la clase \texttt{EntradaLightFM} se tiene:
\begin{itemize}
\tightlist
\item \texttt{leer\_x}: estos métodos recogen los datos de los \textit{.csv} para cada conjunto de datos. Estos métodos son utilizados por los métodos de obtención de matrices de la clase \texttt{SistemaLightFM}.
\end{itemize}

En el archivo \texttt{Salida} se tiene:
\begin{itemize}
\tightlist
\item \texttt{imprimir\_resultados\_clasico}: este método imprime las métricas del modelo clásico escogido.
\end{itemize}

En el archivo \texttt{Persistencia} se tiene:
\begin{itemize}
\tightlist
\item \texttt{guardar\_matrices}: este método guarda en un archivo \textit{pickle} las matrices que \textit{LightFM} necesita para obtener los modelos.
\item \texttt{guardar\_modelo}: este método guarda el modelo obtenido por \textit{LightFM} en un archivo \textit{pickle}.
\end{itemize}

En el archivo \texttt{Interfaz} se tiene:
\begin{itemize}
\tightlist
\item \texttt{elegir\_dataset}: este método muestra un menú mediante el cual elegimos un conjunto de datos que utilizar.
\item \texttt{elegir\_modelo}: este método muestra un menú mediante el cual elegimos un modelo concreto a crear. 
\item \texttt{main\_clasico}: programa principal si el modelo escogido es \textit{LightFM}. 
\end{itemize}

En la clase \texttt{SistemaLightFM} tenemos:
\begin{itemize}
\tightlist
\item \texttt{matrices\_x}: estos métodos crean las matrices necesarias para cada conjunto de datos. 
\item \texttt{modelo\_x}: estos métodos crean los distintos modelos de recomendación. 
\item \texttt{resultados\_x}: estos métodos obtienen los resultados para cada modelo de recomendación. 
\end{itemize}

\subsection{Diseño con Spotlight}\label{diseño-spotlight}
Para la parte del modelo basado en aprendizaje profundo, se tienen los siguientes archivos:
En la clase \texttt{EntradaSpotlight} se tiene:
\begin{itemize}
\tightlist
\item \texttt{leer\_x}: estos métodos recogen los datos de los \textit{.csv} para cada conjunto de datos. Estos métodos son utilizados por los métodos de obtención de interacciones de la clase \texttt{SistemaSpotlight}.
\end{itemize}

En el archivo \texttt{Salida} se tiene:
\begin{itemize}
\tightlist
\item \texttt{imprimir\_resultados\_dl}: este método imprime las métricas del modelo basado en aprendizaje profundo escogido.
\end{itemize}

En el archivo \texttt{Persistencia} se tiene:
\begin{itemize}
\tightlist
\item \texttt{guardar\_interacciones}: este método guarda en un archivo \textit{pickle} las interacciones que \textit{Spotlight} necesita para obtener los modelos.
\item \texttt{guardar\_modelo\_dl}: este método guarda el modelo obtenido por \textit{Spotlight}.
\end{itemize}

En el archivo \texttt{Interfaz} se tiene:
\begin{itemize}
\tightlist 
\item \texttt{main\_dl}: programa principal si el modelo escogido es \textit{Spotlight}. 
\end{itemize}

En la clase \texttt{SistemaSpotlight} tenemos:
\begin{itemize}
\tightlist
\item \texttt{interacciones\_x}: estos métodos crean las interacciones necesarias para cada conjunto de datos. 
\item \texttt{modelo\_x}: estos métodos crean los distintos modelos de recomendación. 
\item \texttt{resultados\_x}: estos métodos obtienen los resultados para cada modelo de recomendación. 
\end{itemize}


\apendice{Documentación técnica de programación}

\section{Introducción}
En este apéndice se explica cómo está estructurado el proyecto en cuanto a los directorios, así como los pasos para compilar, instalar y ejecutar el programa.

\section{Estructura de directorios}
El proyecto está dividido en los siguientes directorios:
\begin{itemize}
\tightlist
\item \textbf{/:} carpeta raíz. Contiene la carpeta \textit{docs}, la carpeta \textit{src} y los ficheros \textit{.gitignore} y \textit{README.md}.
\item \textbf{/docs:} contiene todos los archivos relacionados con la documentación tanto en \LaTeX{} como en \textit{pdf}. También contiene los archivos de bibliografía y las imágenes que se utilizan a lo largo de la documentación.
\item \textbf{/docs/img:} carpeta que contiene las imágenes necesarias para apoyar la documentación.
\item \textbf{/docs/tex:} carpeta que contiene los distintos apartados de la memoria y los anexos en \LaTeX{}.
\item \textbf{/src:} carpeta que contiene todo el código de la aplicación.
\item \textbf{/src/controlador:} contiene los ficheros \textit{.py} con los sistemas de \textit{LightFM} y \textit{Spotlight}.
\item \textbf{/src/modelo:} contiene los ficheros \textit{.py} correspondientes a la \textit{Entrada}, \textit{Salida} y \textit{Persistencia} de los datos.
\item \textbf{/src/vista:} contiene los ficheros \textit{.py} correspondientes a la \textit{Interfaz} (de texto), \textit{Flask} (web) y \textit{Forms} (formularios de los que se compone la interfaz web). Además, contiene los \textit{.html} de la interfaz.
\item \textbf{/src/vista/templates:} contiene las páginas \textit{.html} que componen la interfaz web de la aplicación.
\item \textbf{/src/uploads:} carpeta donde se guardan los modelos y las matrices obtenidas.
\item \textbf{/src/notebooks:} esta carpeta contiene archivos que en su día se utilizaron de prueba. Se podrían borrar, pero se prefiere dejarlos para que haya constancia de que se trabajó en ellos.
\end{itemize}

\section{Manual del programador}
Las herramientas necesarias para la realización del proyecto han sido:
\begin{itemize}
\tightlist
\item Python 3.6
\item Spyder o Sublime Text
\item LightFM
\item Spotlight
\item Flask
\item WTForms
\end{itemize}
A continuación se detalla como instalar estas herramientas:

\subsection{Python}
Lenguaje de programación escogido para llevar a cabo el proyecto. Se puede descargar la última versión desde el siguiente enlace: \href{https://www.python.org/downloads/}{Descargar Python}.

\subsection{IDE}
Se puede utilizar el IDE que más guste. En este caso, se ha utilizado tanto \textit{Spyder} como \textit{Sublime Text}, además de los notebooks de \textit{Jupyter}.

\textit{Spyder} y \textit{Jupyter} vienen dentro del paquete de \textit{Anaconda} (que también incluye Python). \href{https://www.datacamp.com/community/tutorials/installing-anaconda-windows}{Guía para instalar Anaconda en Windows}. En este otro enlace se descarga \textit{Anaconda} directamente para cualquier sistema operativo: \href{https://www.anaconda.com/distribution/}{Descargar Anaconda}.

\textit{Sublime Text} se puede obtener desde el siguiente enlace: \href{https://www.sublimetext.com/3}{Descargar Sublime Text}.

\subsection{LightFM}
Librería utilizada para la obtención de los modelos de recomendación clásicos. Para instalar la librería es recomendable seguir los pasos descritos en el siguiente enlace: \href{http://lyst.github.io/lightfm/docs/home.html#installation}{Instalación de LightFM}.

\subsection{Spotlight}
Librería utilizada para la obtención de los modelos de recomendación basados en aprendizaje profundo. Para instalar \textit{Spotlight} basta con ejecutar el siguiente comando: \textbf{conda install -c maciejkula -c pytorch spotlight=0.1.5} \cite{instalacion-spotlight}. Como se puede observar, con este comando también instalamos \textit{PyTorch}.

\subsection{Flask}
Con \textit{Flask} se obtiene la interfaz web. Tan solo hay que ejecutar el comando \textbf{pip install Flask} para instalarlo \cite{instalacion-flask}.

\subsection{WTForms}
Esta librería se utiliza para poder generar todos los formularios que componen la interfaz web. Para instalarla vale con ejecutar \textbf{pip install WTForms} \cite{instalacion-flask}.

\section{Compilación, instalación y ejecución del proyecto}
\subsection{Obtención del código}
El código se encuentra disponible en el siguiente repositorio de GitHub: \href{https://github.com/rnc0011/SistemaRecomendacionTFG}{Proyecto}. Una vez dentro del repositorio, basta con descargar el proyecto desde el botón de descarga o clonación.

\imagen{repositorio}{Descargar proyecto}

\subsection{Instalación}
Para poder compilar y ejecutar el proyecto es necesario tener instaladas todas las herramientas listadas anteriormente.

\subsection{Compilación y ejecución}
Una vez descargado el repositorio e instaladas todas las herramientas necesarias, se puede proceder a ejecutar la aplicación. Para ello, navegamos desde la consola hasta la carpeta \textit{src} del proyecto:

\imagen{navegacion_proyecto}{Navegación hasta \textit{src}}

Para ejecutar el proyecto basta con ejecutar el archivo \textit{\_\_main\_\_.py} que se encuentra dentro de \textit{src}:

\imagen{ejecución_proyecto}{Ejecución del proyecto}

Metiendo la opción \textit{2} ejecutamos la interfaz web obtenida gracias a \textit{Flask}. Se tiene que abrir el navegador e ir a la dirección: \textbf{http://localhost:5000/home}.

\imagen{interfaz_web}{Interfaz web de la aplicación}

\section{Pruebas del sistema}
-
\apendice{Documentación de usuario}

\section{Introducción}
En este apéndice se explica cómo el usuario debe instalar y utilizar la aplicación. También se indican los requisitos que necesita el usuario para llevarlo a cabo.

\section{Requisitos de usuarios}
\begin{itemize}
\tightlist
\item Navegador web.
\item Conjuntos de datos, en caso de que se quieran utilizar otros además de los de prueba.
\end{itemize}

\section{Instalación}

\section{Manual del usuario}





\bibliographystyle{plain}
\bibliography{bibliografiaAnexos}

\end{document}
